\subsection{Vis�o Geom�trica}

Na se��o anterior demonstramos que o Operador de D'Alembert nas novas coordenadas:
%
%
\[ \square =  (d^2 - c^2)\deli{^2}{{t'}^2} - (a^2 - b^2)\deli{^2}{{x'}^2} + 2(bd - ac)\deli{^2}{x'\partial t'}\]
%
%
O `novo' \textbf{Operador de D'Alembert} deve ser $\square' = \partial^2/\partial {t'}^2 - \partial^2/\partial {x'}^2$ de forma que se tenha $\square \psi(x,t) = \square' \varphi(x',t') = 0$, logo � necess�rio que:
%
%
%
\[bd - ac = 0 \qquad d^2 - c^2 = 1  \qquad a^2 - b^2 = 1\]
%
%
%
Segundo a rela��o $\cosh^2\theta - \sinh^2\theta = 1$ podemos definir que:
%
%
%
\[d^2 - c^2 = 1  \Rightarrow d = \cosh\theta \ , \ c = \sinh\theta\]
%
%
%
Da mesma maneira podemos ainda definir que:
%
%
%
\[a^2 - b^2 = 1  \Rightarrow a = \cosh\phi \ , \ b = \sinh\phi\]
%
%
%
Da rela��o $bd - ac = 0$ pode-se concluir que:
%
%
%
\[\frac{b}{a} = \frac{c}{d} \Rightarrow \tanh\phi = \tanh\theta \Leftrightarrow \phi = \theta\]
%
%
%
Dessa maneira a transforma��o linear (Transforma��o de Lorentz) nada mais � do que uma rota��o no plano:
%
%
%
\[ \begin{pmatrix}  x' \\ y' \end{pmatrix} = \begin{pmatrix} \cosh\theta & \sinh\theta \\ \sinh\theta &  \cosh\theta \end{pmatrix} \begin{pmatrix}  x \\ y \end{pmatrix} \]
%
%
%